\documentclass[a4paper, notitlepage, 11pt]{article}

\usepackage[a4paper, left=17mm, bottom=17mm, right=35mm, top=17mm]{geometry}

\input{includes-common}
\newfontfamily\papercodeblockfont{JetBrainsMono Nerd Font}[Scale=MatchLowercase]
\renewcommand{\codeblockfont}{\papercodeblockfont}

\usepackage[backend=biber, natbib=true, style=apa]{biblatex}
\addbibresource{refs.bib}

\usepackage{float}

\graphicspath{{figures/}}

%%% Paper information
\title{Habit cycles are infodesics:\\ segmentation and switching costs in twisted worlds}
\author{%
  Karen Archer\textsuperscript{1},\,
  [co-authors TBC]%
}
\date{DRAFT SKELETON}

\makeatletter
\renewcommand{\@maketitle}{%
  \begin{center}
    {\Large\bfseries \@title \par}
    \vspace{0.75em}
    {\normalsize \@author \par}
    \vspace{0.5em}
    {\small \textsuperscript{1}Adaptive Systems Group, University of Hertfordshire, Hatfield, UK \par}
    \vspace{0.5em}
    {\small \@date \par}
  \end{center}
  \vspace{1em}
}
\makeatother

% ---------------------------------------------------------------------------
% SKELETON, seeded from docs/habit-cycles-as-infodesics.md (2026-07).
% Paper 1 (twists-home-vectors) reports the habit structures empirically;
% this paper gives the infodesic reading: why the structures are what they
% are.  TODO markers reference the source assets:
%   [NB:policy-switching]  gridFour/notebooks/decision-information/policy-switching.ipynb
%   [NB:infodesics]        gridFour/notebooks/geometry-paper/infodesics/*
%   [GB:01-03]             gridBench/notebooks/goal-geometry/01..03 (drift, corner law,
%                          switching-cost matrix + log-domain fixed-reference solver)
%   [S51]                  gridFour/notebooks/attractor-fingerprint-probe/51 (policy view)
%   [RN:salient]           twists-home-vectors/research-notes.md (doorway salience,
%                          rho = -0.76 with mean_free)
%   [CG]                   cognitive-geom-paper/cognitive-geometry.tex §Infodesics
% ---------------------------------------------------------------------------

\begin{document}

\maketitle

\begin{abstract}
% DRAFT — one paragraph, to be rewritten once results are in.
A twist relabels an agent's actions per state.  Optimised twists build
\emph{habit cycles}: attractors that a single repeated action label drains
the environment into.  A companion paper characterises these structures
empirically; this paper explains them.  We show that a repeated-label flow
is a pure \emph{infodesic} in the space-of-goals geometry --- a route
followable with no state information --- and that the asymmetric homing
mode is a \emph{cost infodesic} segmented at informationally salient
waypoints: the free-energy triangle inequality fails across goal-specific
policies, policy switching becomes advantageous, and the doorway states at
which switches occur are exactly where the evolved label vocabulary changes
hands.  Curvature, concentrated at walls, is the knob: where the world is
flat the pure infodesic, the value geodesic, and the cost infodesic
coincide and homing saturates; where walls concentrate curvature the desics
split, routes segment into single-label legs, and the home-plus-escape
vocabulary of the companion paper emerges.  Switching costs quantify the
segmentation: the cost of abandoning a habit decomposes into an ambient
term set by the environment's drift and a dominant policy-specific term set
by the habit being abandoned.
\end{abstract}

\section{Introduction}
\label{sec:intro}

% TODO: city-directions -> habit framing, one paragraph; then the gap:
% paper 1 reports WHAT evolved twists build (homing labels, coverage
% typology, label roles); this paper asks WHY those structures and not
% others.  The answer is geometric: they are the infodesics of the
% free-energy geometry.
%
% Key sentence from the note: this is not a resurrected side-thread; the
% infodesic is the live, central object of the cognitive-geometry
% programme, continuous from the Space of Goals paper (2022) through the
% current geodesic--loxodrome--curvature work.

\section{Background: the free-energy geometry and its desics}
\label{sec:background}

% TODO [CG]: import the definitions --
%   free energy F as quasi-metric on goals; pure infodesic (state-blind
%   route, I_D -> 0); cost infodesic (F-optimal route, triangle equality);
%   epsilon-infodesic gap; value geodesic; loxodrome/constant-bearing
%   reading.
% TODO [NB:policy-switching]: restate the concluded result -- the
% free-energy triangle inequality FAILS across goal-specific policies and
% multi-policy processing sharply reduces information cost.  This is the
% formal seed of segmentation.  Re-verify on the gridcore kernel (fast;
% the fixed-reference log-domain solver in [GB:03] is the tool).

\section{Single-label flows are pure infodesics}
\label{sec:pure}

% From the note, near-verbatim draft:
A pure infodesic is a sequence visitable with no state information ---
action selection independent of state, decision information
$\ItgD \to 0$: following the same direction with as little cognitive
correction as possible.  A single-label flow is exactly that.  Press one
label in every state: the label is state-independent even though the twist
makes the \emph{physical} direction vary cell to cell.  The home vector of
the companion paper is the direction field of a pure infodesic; the habit
cycle is the attractor the pure infodesic drains into --- its terminus.
The habit is the whole pure-infodesic object: flow plus limit set.

% TODO figure: single-label flow on the exemplar twist, annotated as an
% infodesic (reuse [S51] policy-view panel 1 machinery, goal-free variant).

\section{The asymmetric mode is a segmented cost infodesic}
\label{sec:segmentation}

% From the note -- the role mapping to be expanded into prose + table:
%   home label            <-> one pure-infodesic leg
%   directed-escape label <-> the policy switch to a second leg
%   doorway               <-> informationally salient midpoint / subgoal
%   silenced label        <-> the deficient / trivial part
%
% TODO [S51]: the policy-view routes ARE the segmentation figure -- greedy
% routes switch labels at room seams; each leg is single-label.
% TODO [RN:salient]: doorway-salience correlates with low free energy
% (Spearman rho = -0.758 with mean_free over the salience-evaluated pool)
% -- the waypoints the segmentation predicts are the states selection
% already favours routing through.  Recompute on the clean cohort.

\section{Switching costs quantify the segmentation}
\label{sec:switching}

% TODO [GB:01-03]: port the switching-cost machinery --
%   SC(A->B) = F_B^{ref=pi_A}(A), log-domain fixed-reference solver;
%   H1 ambient drift predicts a minority share (r = 0.52, slope 7.3);
%   H2 policy-specific asymmetry dominates (14.5x the uniform-prior skew);
%   H3 the sink field survives the change of reference (r = 0.68).
% Reading: the cost of abandoning a habit = ambient floor (environment
% drift) + dominant policy-specific term (the particular habit abandoned).
% TODO: notebook 04 (repertoire decomposition I(A;S,G) = I(A;S) + I(A;G|S))
% belongs here if it matures.

\section{Curvature explains the bimodality}
\label{sec:curvature}

% From the note:
% - flat worlds (wrap_grid): pure infodesic = value geodesic = cost
%   infodesic coincide -> one constant-direction label is also shortest ->
%   total saturation, no switching.  The habit IS the geodesic.
% - walled worlds: curvature concentrated at walls/doorways -> the desics
%   split, the cost infodesic goes deficient -> routes segment -> the
%   home + directed-escape structure.  The deficiency IS the need for the
%   escape label.
% - a twist is a global action labelling on a possibly non-flat space; its
%   holonomy is the I_D cost the geometry forces (chi as holonomy: TODO
%   formalise or demote to remark).
% - beta places evolved habits at the loxodromic (cheap, far-predictable)
%   end of the loxodrome <-> geodesic interpolation.

\section{Experiments}
\label{sec:experiments}

% TODO plan (all machinery exists):
% 1. Triangle-inequality violation, systematically: fraction of goal
%    triples violating, by environment and beta [NB:policy-switching,
%    gridcore port].
% 2. Segmentation observed: label-switch points along greedy routes vs
%    doorway/salience sets on the clean cohort [S51 + RN:salient].
% 3. Switching-cost decomposition across the palette [GB:03 scaled up].
% 4. Curvature axis: flat -> pillar -> four_rooms sweep; where do the
%    desics split? (new, but cheap with the gridcore kernel).

\section{Discussion}
\label{sec:discussion}

% TODO: habits as the cheap extreme of the infodesic family; place-graph
% (Mallot) as the segmentation's navigation reading; bridge to QD /
% skill-composition as the search-side continuation.

\printbibliography

\end{document}
